% generated by scripts/paper_supp_speed.py (edit the script, not this file); numbers: docs/paper/supp_speed_numbers.json
\section{Where the speed limit is}
\label{sec:suppl:speed}

Gaussian map fitting is the largest kind of work in run~1: 82.6--84.6\,min per half (\cref{tab:suppl:kinds}).
% src: docs/paper/supp_timing_kinds.csv (map, pro6000_h1/h2); recomputed in docs/paper/supp_speed_numbers.json (runs.pro6000.kinds)
Each update of a KB map fit (cold fit and P3) evaluates the Fourier transform of all $K=48{,}000$ Gaussians (all symmetry copies) at up to 262{,}144 sampled grid frequencies and back-propagates through it:
$1.26\times10^{10}$ Gaussian--frequency pairs forward and again backward.
% src: runs/engine/supp_speed_20261121/profile/p3_base_prof.log (argv --K 48000; 'grid-frequency samples <=262144/update'); stags/train.py (KB objective sample_size=262144); stags/normal_kb.py (objective on model.expanded())
The image batch does not enter this cost: the KB statistics accumulate all images once, before the fit.
% src: stags/train.py fit_half (grid_stats.objective(model, kmax, sample_size, generator), no image batch); stags/normal_kb.py build_kb_statistics

\paragraph{Profile.}
On one PRO 6000 a P3 update takes 35.4\,ms (28.3 it/s): 13.0\,ms in the forward and 18.0\,ms in the backward kernel (97.7\% of the GPU time), 0.7\,ms in other GPU work and 3.6\,ms of GPU idling during host synchronizations (\cref{tab:suppl:speedprofile}).
% src: runs/engine/supp_speed_20261121/profile/p3_base_prof_cuda_top.txt; runs/engine/supp_speed_20261121/profile/p3_base_prof_gpu_timeline.json; runs/engine/supp_speed_20261121/profile/p3_base_prof.log; docs/paper/supp_speed_profile.csv
Per pair, both kernels evaluate an exponential and a hardware sine and cosine, and they read only the Gaussians and the frequencies, so they are arithmetic-bound by design (we did not collect hardware counters).
% src: stags/fourier_volume_triton.py (forward/backward: tl.exp, sin.approx.f32/cos.approx.f32; inputs MU, COEF, MASS, KX/KY/KZ)
Consistent with this, the cold fit of the RTX 4090 run averages 24.4 / 23.9 it/s, close to the 25.6 / 25.5 it/s of the PRO 6000 run (halves 1 / 2, including start-up), whose Max-Q cards are limited to 300\,W.
% src: runs/engine/pose_abc_adopt_20261113/blind/kb/h{1,2}/fit0.log (last 'rate' line); runs/engine/abc_4090_20261117/full_logs/kb/h{1,2}/fit0.log (last 'rate' line); docs/paper/supp_hardware.txt (Max-Q Workstation Edition); runs/engine/supp_speed_20261121/profile/pro6000_power_limit.txt (power.limit)

\paragraph{Exact speed-ups.}
Four implementation changes keep batches, schedules, random draws and update counts: no host synchronization per update in the map-fit loops, 16 instead of 32 Gaussians per backward-kernel program, KB-statistics tiles of 262{,}144 instead of 32{,}768 samples, and a single cast when loading images.
% src: runs/engine/supp_speed_20261121/profile/speed_profile.diff (stags/train.py, stags/engine.py, stags/fourier_volume_triton.py BWD_BLOCKS, stags/normal_kb.py sample_chunk, stags/relion5_io.py)
With frozen KB statistics and initial map, 300 P3 updates are bitwise identical (objective history, parameters, optimizer state, map); the tiled KB build only reorders float32 atomic additions.
% src: runs/engine/supp_speed_20261121/profile/pro_eq_compare.json (history_bitwise_equal True, model bitwise True, optimizer_state_max_abs_diff 0.0, B_4536.mrc bitwise True); runs/engine/supp_speed_20261121/profile/pro_bench_eq.txt (EQ line); the test covers the fit_half (cold fit/P3) path and the backward kernel, not the ladder loops in stags/engine.py
On the PRO 6000 the P3 fit gains 1.6\% and the cold fit 2.2\% in it/s, and the KB build is 2.7--3.1$\times$ faster; we estimate a saving of about 2.2\,min per run (0.9\% of run~1). No reported run used these changes.
% src: runs/engine/supp_speed_20261121/profile/pro_bench_eq.txt (load_stacks not faster: 44.73 vs 47.12 s, not claimed); estimate: docs/paper/supp_speed_numbers.json projected_saving_min (run-1 update counts x (1/rate before - 1/rate after) + KB build difference, one cold fit + four P3 per half; halves in parallel); backward tile on an RTX 5070 Ti: 38.84 -> 34.79 ms, max rel diff 0.0 (runs/engine/supp_speed_20261121/profile/rtx5070ti_bench_kernels2.log)
Without any idling the update would reach 31.5 it/s, and the kernels alone cap it at 32.2 it/s: the implementation is near its kernel-bound floor.
% src: docs/paper/supp_speed_numbers.json (profile.no_idle_rate = 1/GPU busy time, profile.kernel_floor_rate = 1/(forward+backward))
The ${\sim}200$ it/s of real-space 3D Gaussian splatting would need about 7$\times$ fewer pair evaluations; in Fourier space a Gaussian is not compact, so every Gaussian meets every sampled frequency. Fewer sampled frequencies per update, or splatting onto the padded grid with one FFT per update, could provide this; both change the recipe and are untested.
% src: docs/paper/supp_speed_numbers.json (profile.unprofiled_rate); 200 it/s is the rate quoted to us for 3DGS-style fits, not measured here

\paragraph{Fewer updates and fewer Gaussians.}
Two predeclared variants (\cref{tab:suppl:speed}) each completed once, blind from the picks, on the machine of run~1. We compare them with the range of runs~1 and~2, our only estimate of the run-to-run spread (crossings 3.269/3.251 and 3.294/3.245\,\AA).
% src: engine_rules.json steps.supp_speed_20261121 (variants, question, reporting, limitations, incidents); runs/engine/pose_abc_adopt_20261113/blind/scores/cycle4/results.json; runs/engine/abc_4090_20261117/scores/cycle4/results.json
\texttt{--fit\_scale 0.5} halves the updates of the cold fit and of the four P3 fits (ladder unchanged). It saves 10.1\,min against run~1 (246.1\,min); the cold fit and P3 shrink only from 19.4 to 12.4 and from 27.3 to 20.2\,min, because their setup (loading, reconstruction, KB build) is fixed.
% src: docs/paper/supp_speed_variants.csv; report.json wall_min, fit_min and cycles_minutes[*].p3 (summed) (runs/engine/supp_speed_20261121/fit_scale_0.5/report.json; runs/engine/pose_abc_adopt_20261113/blind/report.json); updates {1: 12096, 2: 12016} / 3024 vs {1: 24192, 2: 24032} / 6048; cycle-4 P3 half 1, 'Elapsed' vs. update-loop 'elapsed': run 1 377 vs. 253 s, fit_scale 266 vs. 140 s (runs/engine/pose_abc_adopt_20261113/blind/kb/cycle4/p3_h1.log; runs/engine/supp_speed_20261121/fit_scale_0.5/kb/cycle4/p3_h1.log)
Of its seven final-map values, one lies inside the range of runs~1 and~2, five are better than both runs (by at most 0.014\,\AA{} and 0.008 in FSC) and one is worse than both runs (half-2 FSC at 3.5\,\AA, by 0.001 in FSC).
% src: docs/paper/supp_speed_numbers.json (spread_check.fit05: value vs. [min, max] of runs 1 and 2); runs/engine/supp_speed_20261121/fit_scale_0.5/scores/cycle4/results.json
\texttt{--k\_scale 0.5} halves every Gaussian budget: the ladder's densification targets and caps, and $K$ from 48{,}000 to 24{,}000 in the cold fit and P3.
% src: runs/engine/supp_speed_20261121/k_scale_0.5/launch.json (recipe_scales); engine_rules.json steps.supp_speed_20261121.variants.k_half
Fewer Gaussians do not make the ladder cheaper: a seed-stage M-step round takes 62--66\,s with $K=528$, as with $K=1{,}056$ in run~1 (65--66\,s), but half~1 needs 24 instead of 9 rounds (28.2--29.7 instead of 11.9\,min) to meet the seed stage's stop rule (half~2: 10 rounds).
% src: 'round N timing: .. M-step' lines and K in lean/logs/seed_h{1,2}.log, lean/timing.json stages.step1_seed.halves_min of runs/engine/pose_abc_adopt_20261113/blind/lean and of runs/engine/supp_speed_20261121/k_scale_0.5.aborted-gpu-contention/lean; runs/engine/supp_speed_20261121/k_scale_0.5.aborted-stage-budget/lean; stop rule: lattice content 0.12 (engine_rules.json steps.supp_speed_20261121.incidents correction); docs/paper/supp_speed_numbers.json (k_scale_seed)
In both of two attempts, the pipeline's stage-time gate (1.5$\times$ the stage budget) then aborted the next ladder stage, because half~2's stage time includes its wait for half~1. The run reported here is a third attempt with wider stage budgets and an otherwise identical recipe (the gate only guards against stuck runs).
% src: runs/engine/supp_speed_20261121/k_scale_0.5.aborted-gpu-contention/lean/chain.log (step2_16_12: half 2 stage time 22.5 min > 1.5 x budget 15 min); runs/engine/supp_speed_20261121/k_scale_0.5.aborted-stage-budget/lean/chain.log (step2_16_12: half 2 stage time 22.6 min > 1.5 x budget 15 min); runs/engine/supp_speed_20261121/k_scale_0.5/launch.json (argv --ladder_budget step1_seed=45, step2_16_12=40, step2_10=30, step3_8_6.5=60, step4_5.5=80); engine_rules.json steps.supp_speed_20261121.incidents (correction: not GPU contention)
The third attempt was still in progress when this section was generated; its row in \cref{tab:suppl:speed} is pending.
% src: runs/engine/supp_speed_20261121/k_scale_0.5/run_abc-k05.log (no 'pipeline DONE' line yet)
With one completed run per variant, differences of the size of the run-to-run spread cannot be attributed to the variant, and no variant changes the recipe of the main paper.\footnote{In the \texttt{--fit\_scale} run the scorer refused to start without its self-test; after the run the self-test passed and all stages were scored with the same scorer.}
% src: runs/engine/supp_speed_20261121/fit_scale_0.5/scores/cycle4.out (A4 self-test required); runs/engine/supp_speed_20261121/queue.log; runs/engine/supp_speed_20261121/a4_selftest_20261122.log

\begin{table}[t]
\centering
\footnotesize
\setlength{\tabcolsep}{4pt}
% generated by scripts/paper_supp_speed.py; csv: docs/paper/supp_speed_profile.csv
\begin{tabular}{@{}lrr@{}}
\toprule
One P3 map-fit update, one PRO 6000 & ms & share (\%) \\
\midrule
% src: runs/engine/supp_speed_20261121/profile/p3_base_prof_cuda_top.txt (_Volume forward / _VolumeBackward, 20 calls); runs/engine/supp_speed_20261121/profile/p3_base_prof_gpu_timeline.json (union of GPU activity in the trace); step = 1/rate of updates 20..120 in runs/engine/supp_speed_20261121/profile/p3_base_prof.log ('[speed] .. iterations 20..120'), before the profiled updates 300..320
Forward kernel & 13.01 & 36.8 \\
Backward kernel & 18.02 & 51.0 \\
Other GPU work (optimizer, small ops) & 0.72 & 2.0 \\
GPU idle (host synchronization) & 3.61 & 10.2 \\
\midrule
Update (28.28 it/s, not profiled) & 35.36 & 100.0 \\
\bottomrule
\end{tabular}
\\[3pt]
\begin{tabular}{@{}lrrr@{}}
\toprule
Benchmark, one PRO 6000 & before & after & gain \\
\midrule
% src: runs/engine/supp_speed_20261121/profile/pro_bench_eq.txt (iterations 100..1200; build_kb_statistics; load_stacks)
P3 map fit (it/s) & 27.98 & 28.43 & +1.6\% \\
Cold map fit (it/s) & 27.94 & 28.56 & +2.2\% \\
KB statistics, P3 (s) & 30.5 & 9.8 & 3.1$\times$ \\
KB statistics, cold fit (s) & 27.1 & 10.2 & 2.7$\times$ \\
\bottomrule
\end{tabular}
\caption{\textbf{Map-fit profile.} Top: one update of the P3 warm map fit (torch.profiler over 20 updates; shares of the un-profiled update time). Bottom: steady-state rate (updates 100--1200) and KB statistics build time before and after the exact speed-ups of \cref{sec:suppl:speed} (full half-1 data, recorded commands).}
\label{tab:suppl:speedprofile}
\end{table}

\begin{table*}[t]
\centering
\footnotesize
\setlength{\tabcolsep}{3.5pt}
\providecommand{\SpeedPending}{\textit{pending}}
% generated by scripts/paper_supp_speed.py; csv: docs/paper/supp_speed_variants.csv
\resizebox{\textwidth}{!}{%
\begin{tabular}{@{}llrcrcccccc@{}}
\toprule
& GPUs & $K$ & Updates, cold / P3 & Wall (min) & Map fit (min) & FSC 0.143 (\AA) $\downarrow$ & FSC 3.5\,\AA\ $\uparrow$ & Band FSC $\uparrow$ & Atomic 0.5 (\AA) $\downarrow$ & Atomic 0.5, ${<}14$\,\AA\ (\AA) $\downarrow$ \\
\midrule
% src: K/updates runs/engine/pose_abc_adopt_20261113/blind/kb/h{1,2}/fit0.log, cycle1/p3_h1.log (argv); wall/kinds runs/engine/pose_abc_adopt_20261113/blind/report.json; map fitting = kinds() of scripts/paper_supp_timing.py over runs/engine/pose_abc_adopt_20261113/blind/lean/logs/{seed,ladder}_h{1,2}.log; runs/engine/pose_abc_adopt_20261113/blind/kb/h{1,2}/fit0.log; cycle*/p2_h*/search.json + p2_h*.log; scan_h*/optics.json + scan_h*.log; p1_h*/p1.json; p3_h*.log; scores runs/engine/pose_abc_adopt_20261113/blind/scores/cycle4/results.json + structure/curves.csv; atomic finer than 14 A: atomic_detail() over the atomic_avg curve
Run 1 & PRO 6000 & 48{,}000 & 24{,}032--24{,}192 / 6{,}048 & 256.2 & 84.6 / 82.6 & 3.269 / 3.251 & 0.247 / 0.243 & 0.228 / 0.226 & 4.04 & 4.04 \\
% src: K/updates runs/engine/abc_4090_20261117/full_logs/kb/h{1,2}/fit0.log, cycle1/p3_h1.log (argv); wall/kinds runs/engine/abc_4090_20261117/full_logs/report.json; map fitting = kinds() of scripts/paper_supp_timing.py over runs/engine/abc_4090_20261117/full_logs/lean/logs/{seed,ladder}_h{1,2}.log; runs/engine/abc_4090_20261117/full_logs/kb/h{1,2}/fit0.log; cycle*/p2_h*/search.json + p2_h*.log; scan_h*/optics.json + scan_h*.log; p1_h*/p1.json; p3_h*.log; scores runs/engine/abc_4090_20261117/scores/cycle4/results.json + structure/curves.csv; atomic finer than 14 A: atomic_detail() over the atomic_avg curve
Run 2 & RTX 4090 & 48{,}000 & 24{,}032--24{,}192 / 6{,}048 & 263.2 & 72.6 / 74.0 & 3.294 / 3.245 & 0.252 / 0.266 & 0.228 / 0.234 & 4.01 & 4.01 \\
% src: K/updates runs/engine/supp_speed_20261121/fit_scale_0.5/kb/h{1,2}/fit0.log, cycle1/p3_h1.log (argv); wall/kinds runs/engine/supp_speed_20261121/fit_scale_0.5/report.json; map fitting = kinds() of scripts/paper_supp_timing.py over runs/engine/supp_speed_20261121/fit_scale_0.5/lean/logs/{seed,ladder}_h{1,2}.log; runs/engine/supp_speed_20261121/fit_scale_0.5/kb/h{1,2}/fit0.log; cycle*/p2_h*/search.json + p2_h*.log; scan_h*/optics.json + scan_h*.log; p1_h*/p1.json; p3_h*.log; scores runs/engine/supp_speed_20261121/fit_scale_0.5/scores/cycle4/results.json + structure/curves.csv; atomic finer than 14 A: atomic_detail() over the atomic_avg curve
% scored after the run with the same scorer: runs/engine/supp_speed_20261121/queue.log ('rescored cycle4 exit 0'); self-test runs/engine/supp_speed_20261121/a4_selftest_20261122.log
\texttt{-{}-fit\_scale 0.5} & PRO 6000 & 48{,}000 & 12{,}016--12{,}096 / 3{,}024 & 246.1 & 69.9 / 67.2 & 3.294 / 3.240 & 0.260 / 0.242 & 0.230 / 0.235 & 4.00 & 4.00 \\
% src: K/updates runs/engine/supp_speed_20261121/k_scale_0.5/kb/h{1,2}/fit0.log, cycle1/p3_h1.log (argv); wall/kinds runs/engine/supp_speed_20261121/k_scale_0.5/report.json; map fitting = kinds() of scripts/paper_supp_timing.py over runs/engine/supp_speed_20261121/k_scale_0.5/lean/logs/{seed,ladder}_h{1,2}.log; runs/engine/supp_speed_20261121/k_scale_0.5/kb/h{1,2}/fit0.log; cycle*/p2_h*/search.json + p2_h*.log; scan_h*/optics.json + scan_h*.log; p1_h*/p1.json; p3_h*.log; scores runs/engine/supp_speed_20261121/k_scale_0.5/scores/cycle4/results.json + structure/curves.csv; atomic finer than 14 A: atomic_detail() over the atomic_avg curve
\texttt{-{}-k\_scale 0.5} & PRO 6000 & 24{,}000 & 24{,}032--24{,}192 / 6{,}048 & 230.0 & 76.8 / 61.1 & 3.235 / 3.280 & 0.222 / 0.218 & 0.223 / 0.206 & 16.13 & 4.09 \\
\bottomrule
\end{tabular}}
\caption{\textbf{Speed-study variants} (one completed run each) against runs 1 and 2 of the adopted recipe. $K$: Gaussians of the cold KB fit and of P3 (all symmetry copies); \texttt{-{}-k\_scale} also halves the ladder's Gaussian budgets. Map fit: minutes per half (half 1 / half 2) of all map fitting (ladder M-steps, cold fit, P3; as in \cref{tab:suppl:kinds}). Scores: final completed cycle, half 1 / half 2 against the opposite RELION-5 half map; band FSC: mean FSC over 3.6--3.28\,\AA; atomic 0.5: first FSC 0.5 crossing of the half-map average against 5L93, and the first crossing between shells finer than 14\,\AA, beyond the band near 15\,\AA\ where the atomic curve of every map dips.}
\label{tab:suppl:speed}
\end{table*}


